\section{Methods}

This section summarizes the main methodological choices used in the analyses. Full mathematical definitions, simulation settings, priors, data-processing steps, and implementation details are provided in the \textit{Supplementary Information}.

\noindent\textbf{Modeling phylodynamic rates with BELLA.} BELLA models selected phylodynamic rates as functions of external predictors using feedforward neural networks embedded in a Bayesian phylodynamic model. In the applications considered here, the underlying process is a birth--death model whose rates may describe epidemiological transmission, sampling, and migration, or macroevolutionary speciation, extinction, and fossil sampling \citep{stadler2013uncovering,Vaughan2025BDMM,Heath2014FBD}. The neural network maps predictors such as time, traits, or mobility covariates to the target rates, which in this study include speciation, extinction, and migration rates. Non-target rates, such as sampling rates and the become-uninfectious rate in epidemiological models, are not modeled by the network and are instead fixed or sampled from their usual priors.

The network weights and biases are assigned priors and sampled jointly with the remaining phylogenetic and phylodynamic parameters using Markov chain Monte Carlo (MCMC). At each MCMC step, BELLA transforms the current predictor values into phylodynamic rates, and these rates are used to evaluate the phylodynamic likelihood. Output activation functions were chosen to keep predicted parameters within their natural support, for example by enforcing positivity or finite bounds. This framework retains standard Bayesian uncertainty propagation while allowing nonlinear and interacting predictor effects that are not available under ordinary linear predictor models.

\noindent\textbf{Simulation study.} We benchmarked BELLA on eight simulated scenarios spanning epidemiological and macroevolutionary settings (Supplementary Information, Section~\ref{subsec:sim-study}). These scenarios tested whether the model could recover constant, linear, nonlinear, and interacting relationships between predictors and phylodynamic rates. The epidemiological simulations included time-varying reproduction numbers and pairwise migration among structured populations, while the fossilized birth--death simulations included time-varying and trait-dependent speciation and extinction rates. Each scenario also included irrelevant random predictors to assess robustness to overfitting.

For every scenario, we simulated 100 trees with 200--500 tips and analyzed each replicate using three classes of inference models: a predictor-agnostic model, generalized linear models (GLMs), and multiple BELLA configurations. The BELLA configurations varied in network size, activation functions, and priors on the weights, allowing us to assess sensitivity to architecture while comparing against simpler and more constrained alternatives. Detailed scenario definitions, parameter values, priors, and MCMC settings are reported in the \textit{Supplementary Information}.

\noindent\textbf{Evaluation metrics.} Model performance in the simulation study was evaluated from posterior summaries of the target phylodynamic parameters. We used mean absolute percentage error (MAPE) to measure estimation accuracy, the coefficient of variation (CV) to quantify relative posterior dispersion, and empirical coverage of 95\% credible intervals to assess uncertainty calibration. To compare computational efficiency, we also computed the mean effective sample size per hour across target parameters and simulation replicates. Formal definitions of these metrics are given in the \textit{Supplementary Information}.

\noindent\textbf{Model interpretation.} We used partial dependence plots (PDPs) and Shapley feature importance to interpret the predictor effects learned by BELLA \citep{friedman2001pdps,lundberg2017unified}. PDPs summarize the marginal relationship between a predictor and the network output, whereas Shapley values quantify the contribution of each predictor to the model predictions. These complementary summaries were used to distinguish relevant predictors from irrelevant covariates and to visualize nonlinear effects recovered by the neural network.

Because BELLA produces a posterior distribution over neural networks rather than a single fitted model, interpretation was performed across posterior samples. For each run, we subsampled neural network configurations from the MCMC output, computed PDPs and Shapley values for each configuration, and summarized the resulting distributions by their medians. The full interpretation procedure is described in the \textit{Supplementary Information}.

\noindent\textbf{Empirical analyses.} We applied BELLA to two empirical datasets. First, we analyzed diversification in New World monkeys using 100 dated phylogenetic trees comprising 121 platyrrhine species, including both extant and extinct taxa \citep{silvestro2019early}. In this analysis, BELLA modeled speciation and extinction rates as functions of geological time and discretized body-mass categories, while sampling rates through time and transition rates among body-mass categories were estimated within the same birth--death framework.

Second, we analyzed the early spread of SARS-CoV-2 in Europe using 123 viral genomes sampled up to 8 March 2020. The phylogeographic analysis used a multi-type birth--death model in which demes represented China, Germany, Italy, France, and other European countries. We compared a GLM and BELLA for modeling migration rates as functions of air-passenger volume normalized by source-country population size, while jointly estimating the phylogeny, reproduction numbers, sampling proportions, and migration dynamics. Posterior migration histories were summarized using phylogenetic particle filtering as implemented in BDMM-Prime \citep{vaughan2019estimating,Vaughan2025BDMM}. Dataset construction, model settings, and convergence criteria for both empirical analyses are provided in the \textit{Supplementary Information}.

\noindent\textbf{Computational runtime.} The increased parameter space associated with neural networks compared to predictor-agnostic or GLM-based formulations makes BELLA computationally slower (Tab.~\ref{tab:ess_per_hour}); in our analyses, the runtime difference varied by roughly a factor of 10.
Smaller neural networks generally converged faster in simple settings, whereas larger architectures were more useful when several predictors or interacting effects were present.
Both empirical analyses were run for 24 hours, by which point all chains had reached a substantial number of effective samples ($\text{ESS} >> 200$) for each estimated parameter.

\noindent\textbf{Implementation and data availability.} BELLA is implemented in the BEAST 2 framework \citep{bouckaert2019beast} as the \textit{Bayesian Evolutionary Layered Learning Architectures} package, available at \url{https://github.com/gabriele-marino/BELLA/}.
The package can be combined with BEAST 2 birth--death population models, sequence evolution models, and MCMC operators, enabling inference either jointly with phylogenetic tree estimation or on fixed, time-calibrated trees.

All code required to reproduce the analyses presented here is available at \url{https://github.com/gabriele-marino/BELLA-companion}. This includes the simulation code, empirical datasets, BEAST configuration files, and post-processing pipelines used for performance evaluation and model interpretation.
